% 日本音響学会 講演論文形式（A4・2段組・10pt）の原稿
% ビルド: cd results/paper_asj && tectonic main.tex   （Tectonic = XeTeX。必要なパッケージは初回に自動で取得）
% 数値の出典: results/GENERATIVE_RESULTS_2026-09-30.md §1・§6・§7、図は scripts/83_fig_fusion_relative.py
% 注意: 学会の正式テンプレートではなく、余白（上下18 mm・左右22 mm・段間6 mm）を合わせた自作の書式
\documentclass[xelatex,ja=standard,a4paper,twocolumn,10pt]{bxjsarticle}
\geometry{top=18mm,bottom=18mm,left=22mm,right=22mm,columnsep=6mm}
\setjamainfont[BoldFont={Hiragino Sans W6}]{Hiragino Mincho ProN}
\setjasansfont[BoldFont={Hiragino Sans W6}]{Hiragino Sans W3}
\usepackage{amsmath,amssymb}
\usepackage{graphicx}
\usepackage{booktabs}
\usepackage[font=small,skip=3pt]{caption}
\usepackage{cite}

\pagestyle{empty}
\setlength{\textfloatsep}{8pt plus 2pt minus 2pt}
\setlength{\floatsep}{6pt plus 2pt minus 2pt}

% 見出しを講演論文らしく小さめに
\makeatletter
\renewcommand{\section}{\@startsection{section}{1}{\z@}{1.0ex plus .3ex}{.5ex}{\normalfont\large\sffamily\bfseries}}
\renewcommand{\subsection}{\@startsection{subsection}{2}{\z@}{.7ex plus .2ex}{.3ex}{\normalfont\normalsize\sffamily\bfseries}}
\makeatother

\newcommand{\figdir}{../figures}

\begin{document}

\twocolumn[{%
  \begin{center}
    {\Large\sffamily\bfseries Whisper系音声認識に向けた喉マイク音声強調における\\[2pt]
      事前学習済み生成モデルとASR損失の役割\footnotemark}\\[8pt]
    {春日　裕次（明治大学　総合数理学部　先端メディアサイエンス学科）}
  \end{center}
  \vspace{6pt}
}]
\footnotetext{Roles of a pretrained generative model and ASR-loss enhancement in throat-microphone speech enhancement for Whisper-family recognizers, by KASUGA, Yuji (Meiji University).}

%────────────────────────────────────────
\section{はじめに}
喉マイクは頸部の振動を接触型センサで捉えるため環境騒音に強いが、軟組織の伝搬特性により高域が大きく失われる。
喉マイクと気導マイクの同時収録コーパスTAPS\cite{taps}では喉マイクが8~kHzで収録されており、4~kHz以上は収録時点から存在しない。
Whisper-small\cite{whisper}で喉マイク音声をそのまま認識すると文字誤り率（CER）は0.446に達する。

前処理としての音声強調（SE）は有力な対策であるが、知覚品質の改善が音声認識（ASR）の改善を保証しないことが知られている\cite{ochiai2024}。
この乖離に対し、凍結したASRの損失でSEを学習する方法が提案されてきた\cite{dissen2024,sato2025}。
しかし、損失に用いた認識器以外への効果は十分に調べられていない。
本稿の実験でも、Whisper-smallの交差エントロピー（CE）損失で学習した喉マイクSEは、WhisperではCERを下げる一方、入力特徴量の異なる認識器では悪化した（4.1節）。
ASR損失で学習したSEは、損失に用いた認識器の入力特徴量に特化しうる。

一方、喉マイクの問題は雑音の混入ではなく帯域や子音の欠落であり、欠けた成分を補う必要がある。
大規模音声で事前学習した生成モデル\cite{ku2024}は、入力が持たない成分を自然な音声として生成できる。
体内伝導マイクではこの種の生成モデルの適応が有効との報告がある\cite{hauret2025rt}が、評価は単一の認識器に限られる。
また、生成的な帯域拡張や声質変換が喉マイク音声のCERを悪化させたという報告もある\cite{yamanaka2026}。

本稿では対象をWhisperと同じ対数メルスペクトログラムを入力とする認識器（以下Whisper系）に定め、損失と設定の選択にはWhisper-smallのみを用い、他のWhisper系5種を評価専用とする。
そのうえで次の3点を示す。
(1) 事前学習済みのflow matching生成モデルをTAPSで適応すると、入力特徴量の異なる認識器を含む9種すべてでCERが下がる。事前学習なしでは悪化し、効果は事前学習に由来する。
(2) 生成モデルの出力（複数サンプルの振幅平均）にWhisper用SEの出力を振幅で融合すると、学習に用いていないWhisper系にも上乗せが得られ、既存のASR損失によるSE（SSL-MSE、Dissen型）より良い。一方、生成モデル自体をWhisperの損失で学習する方法は効かない。
(3) 融合の上乗せは4~kHz未満・以上の両帯域のスペクトル包絡から生じ、他系統の認識器への害は、喉マイクに元々存在しない4~kHz以上の成分から生じる。この害は、Whisper用SEの重みの内挿と対数振幅での融合により、Whisper系の上乗せを保ったままほぼ解消できる。

%────────────────────────────────────────
\section{方法}
\subsection{生成モデルによる復元}
NVIDIAが公開しているflow matching\cite{lipman2023}型の音声復元モデル（4.3億パラメータ、複素STFT領域）\cite{ku2024}を用いる。大規模音声でマスク区間の復元により事前学習されたモデルである。
本研究ではTAPSの喉マイク音声を条件、同時収録の気導マイク音声を目標として追加学習する（ASRの損失は用いない）。
推論では標準正規分布の雑音から出発し、喉マイク音声を条件として常微分方程式をEuler法20ステップで解く。

生成は出発点の雑音に依存して発話ごとにばらつく。
そこで乱数の種を変えた$N$個のサンプルのSTFT振幅を平均し、位相は1つ目のサンプルのものを用いる。
\begin{equation}
  |G| = \frac{1}{N}\sum_{n=1}^{N}|G_n|, \qquad \angle G = \angle G_1
\end{equation}

\subsection{Whisper用SE}
TAPS論文\cite{taps}のSE-Conformer（1210万パラメータ、以下TAPS SE）を初期値とし、次の損失で追加学習したモデル（以下CE-SE）を用いる。
\begin{equation}
  L = L_\mathrm{recon}(\hat{s}, s_\mathrm{air}) + \lambda\, \mathrm{CE}(\mathrm{Whisper}(\hat{s}), y)
\end{equation}
$\hat{s}$はSE出力、$s_\mathrm{air}$は気導音声、$y$は正解テキスト、$L_\mathrm{recon}$はL1波形損失と多重解像度STFT損失の和である。
Whisper-smallは凍結し、$\lambda$はdevのCERで$\{0, 0.1, 0.5, 1, 2, 5, 10\}$から選んで$\lambda=10$とした。

\subsection{振幅融合}
生成モデルの出力$G$（$N=4$の平均）とCE-SEの出力$C$を、STFT（512点、シフト128点）の振幅で重み付き平均し、位相は生成モデルのものを用いる。
\begin{equation}
  Y = \bigl\{(1-w)|G| + w|C|\bigr\}\, e^{j\angle G}
\end{equation}
$w$は、生成1回の出力との融合をdevで比べ、Whisper-smallのCERのみを見て$\{0.25, 0.5, 0.75\}$から選んだ（CERはそれぞれ0.202・0.196・0.198）。
test では同じ$w=0.5$を4サンプル平均との融合に用いた。
融合に新たな学習は要らない。

%────────────────────────────────────────
\section{実験条件}
\subsection{データと学習}
TAPS\cite{taps}の話者独立な分割（train 40話者・dev 10話者・test 10話者、各話者100発話）を用いた。
生成モデルはtrainの4,000対で20k step追加学習した（batch 8、6.14秒の切り出し、学習率$10^{-4}$、EMA、32ビット浮動小数点）。
モデルの選択には最終ステップを用い、testを見て選んでいない。
CE-SEはAdam（学習率$3\times10^{-4}$、batch 4）で最大50エポック学習し、devの損失でearly stoppingを行った。15秒を超える発話は学習から除外した。

\subsection{比較手法}
TAPS SE、CE-SEに加え、2つの既存手法を比較した。
SSL-MSEは、SE出力と気導音声のWavLM-Large\cite{wavlm}全層の表現の二乗誤差を再構成損失に加えてTAPS SEを追加学習したもので、Satoら\cite{sato2025}の方式に従う。
Dissen型は、凍結したASRの損失で前段を学習する方法\cite{dissen2024}に倣い、喉マイクの対数メルスペクトログラムを入力・出力とするU-Net（約700万パラメータ）を、Whisper-smallのCEと気導音声の対数メルへのL1損失で学習したものである。

\subsection{認識器と評価}
\textbf{Whisper系。}
Whisper-small（損失と選択に使用）、Whisper-base、Whisper-medium、Whisper-large-v3-turbo、TAPSの喉マイク音声でファインチューニングしたWhisper-small（FT Whisper）、Qwen3-ASR-1.7Bの6種である。
Qwen3-ASRのエンコーダはWhisperと異なるが、入力は同形式の対数メル（128次元）である。
Whisperはfaster-whisper（CTranslate2）で言語を韓国語に指定して認識した。

\textbf{入力特徴量の異なる認識器。}
Whisper系への特化を確認する対照として、波形入力のCTC型であるMMS-1B\cite{mms}とXLS-R\cite{xlsr}（韓国語で追加学習したもの）、Kaldi型のフィルタバンクを入力とするZipformer\cite{zipformer}を用いた。

\textbf{指標。}
test 1,000発話で、正解と認識結果の双方から句読点を除いたCERを発話ごとに求め（1.0で打ち切り）、平均した。
Zipformerは空白を含めたCERである。
同一話者の発話は独立でないため、話者ごとの平均CERを単位とするWilcoxon符号順位検定（$n=10$、最小の$p=0.002$）を用いた。

%────────────────────────────────────────
\section{実験結果}
\subsection{主な結果}
表\ref{tab:main}に各手法のCERを、図\ref{fig:rel}にTAPS SE比の相対変化を示す。
生成モデル（1回の生成）は9種の認識器すべてでTAPS SEよりCERを15〜24\%下げ、いずれも10話者全員で改善した（$p=0.002$）。
Whisper-smallの認識結果では崩れた発話（CER$\geq$0.99）がなく、出力長も変わらず、置換と挿入の誤りが減っており、生成による作り話の兆候は見られなかった。

CE-SE単体は、Whisper系ではTAPS SEよりCERを下げた（Whisper-baseを除き$-14$〜$-17$\%、baseは$-2.2$\%）。
一方、XLS-Rでは25.8\%、Zipformerでは45.9\%悪化した。
ASR損失による改善は、損失に用いた認識器と入力特徴量の異なる認識器には及ばない。

同じ構造・同じ学習量で初期値を乱数としたモデルは、Whisper-smallで0.551、XLS-Rで0.524とTAPS SEより大幅に悪化した。
生成モデルの効果はモデルの規模ではなく事前学習に由来する。

\begin{table}[t]
  \centering
  \caption{各手法のtest CER。TAPS: TAPS SE、SSL: SSL-MSE、生成×$N$: 生成モデルの$N$サンプル平均、W-FT: FT Whisper、*: 損失・選択に使用}
  \label{tab:main}
  \footnotesize
  \setlength{\tabcolsep}{2.4pt}
  \begin{tabular}{lcccccc}
    \toprule
    認識器 & TAPS & CE-SE & SSL & 生成×1 & 生成×4 & 融合 \\
    \midrule
    W-small* & 0.230 & 0.196 & 0.204 & 0.180 & 0.169 & \textbf{0.163} \\
    W-base & 0.273 & 0.267 & 0.241 & 0.218 & 0.205 & \textbf{0.201} \\
    W-medium & 0.196 & 0.169 & 0.177 & 0.159 & 0.151 & \textbf{0.144} \\
    W-turbo & 0.170 & 0.142 & 0.153 & 0.137 & 0.131 & \textbf{0.125} \\
    W-FT & 0.146 & 0.121 & 0.124 & 0.110 & 0.100 & \textbf{0.092} \\
    Qwen3 & 0.136 & 0.114 & 0.114 & 0.104 & 0.097 & \textbf{0.094} \\
    \midrule
    XLS-R & 0.222 & 0.279 & 0.190 & 0.169 & \textbf{0.161} & 0.163 \\
    MMS-1B & 0.350 & 0.354 & 0.321 & 0.297 & \textbf{0.284} & 0.285 \\
    Zipformer & 0.542 & 0.791 & 0.486 & 0.451 & \textbf{0.433} & 0.483 \\
    \bottomrule
  \end{tabular}
\end{table}

\begin{figure}[t]
  \centering
  \includegraphics[width=0.80\columnwidth]{\figdir/fusion_relative.png}
  \caption{TAPS SE比のCER変化（破線より上がWhisper系）}
  \label{fig:rel}
\end{figure}

\subsection{サンプル平均・融合と既存手法との比較}
4サンプルの振幅平均は、1回の生成よりさらにCERを下げた（Whisper-small 0.180→0.169、10/10話者、$p=0.002$）。
平均CERは乱数の種によらずほぼ一定であったが、発話単位では3回のCERの差が0.1を超える発話が82/1,000あり、平均がこのばらつきを抑えたと考えられる。

CE-SEとの融合は、Whisper系の6種すべてで4サンプル平均よりCERを下げた（$-1.8$〜$-8.0$\%）。
Whisper-medium・large-v3-turbo・FT Whisperでは10/10話者（$p=0.002$）、Whisper-smallで9/10話者（$p=0.004$）で改善し、Whisper-baseとQwen3-ASRでは有意でなかった。
TAPS SE比ではWhisper系で$-26$〜$-37$\%となり、既存手法であるSSL-MSEやCE-SE単体よりも6種すべてで良い。
一方、入力特徴量の異なる認識器では融合の上乗せはなく、Zipformerでは4サンプル平均より11.7\%悪化した（0/10話者）。
融合はWhisper系に特化した上乗せである。

もう1つの既存手法であるDissen型は80次元メルのWhisperにのみ適用できるため、表1とは別に、全条件をtransformers版Whisperの貪欲復号で認識し直して比較した。
Whisper-smallでTAPS SE 0.237、Dissen型 0.239、融合 0.169であり、Dissen型はWhisper-base・medium・FT WhisperでもTAPS SEより悪かった。
喉マイクのメルのみを入力とするため、欠けた帯域を補えなかったと考えられる。

\subsection{Whisperに合わせる成分を学習で組み込む試み}
融合の代わりに、Whisper-smallの損失を学習で組み込む方法を試した（表\ref{tab:fail}）。
まず生成モデル自体を学習した。
flow matchingでは、雑音（$t=0$）から音声（$t=1$）へ向かう経路上の点$x_t$でモデルが速度$\hat{v}$を推定する。
そこから1段で得られる完成音の見積もり$\hat{x}_1 = x_t + (1-t)\hat{v}$にCEをかけて追加学習すると、見積もりのCEは下がるが、20段で生成した出力のCERは悪化した。
CEなしで同じ手順をとると悪化しないため、原因はCE損失にある。
最終段のみに損失をかけるDRaFT-K\cite{draft}型の学習、生成出力に小さな補正モジュールを加える方法も効果がなかった。
次に、生成モデルは固定し、CE-SEを融合後の音に対するCEで再学習した。
Whisper系では一貫して改善したが、幅は0.2〜1.9\%にとどまり、学習なしの融合からの上乗せは小さい。

\begin{table}[t]
  \centering
  \caption{Whisperの損失を学習で組み込む試み（W-small、test CER）}
  \label{tab:fail}
  \footnotesize
  \setlength{\tabcolsep}{3pt}
  \begin{tabular}{lc}
    \toprule
    方法 & 結果 \\
    \midrule
    生成×1（基準） & 0.180 \\
    生成：$\hat{x}_1$にCE（重み0.1／0.01） & 0.193／0.191 \\
    生成：同じ手順でCEなし & 0.179 \\
    生成：DRaFT-K（最終段にCE） & devのCEが悪化 \\
    生成出力に補正モジュール（0.03M） & devのCE変化なし \\
    CE-SEを融合後の音で再学習 & 0.162（融合0.163） \\
    \bottomrule
  \end{tabular}
\end{table}

\subsection{融合の分解}
CE-SEの何がWhisper系に効くのかを、融合の作り方を変えて調べた（表\ref{tab:decomp}）。
CE-SEの出力には、3〜7.75~kHzの250~Hz間隔の位置に鋭いピークが並ぶ櫛状の成分が現れる。このピークをノッチで除いてから融合しても結果はほぼ変わらず、上乗せの原因ではなかった。
4~kHz未満のみ、または4~kHz以上のみを融合すると、Whisper系の上乗せはそれぞれ約半分になり、両帯域から生じていた。
さらに対数振幅をケプストラムの低次（30次未満、スペクトル包絡）と高次（細かい調波構造）に分けて片方だけを融合すると、包絡だけでWhisper系の上乗せの大半が得られ（例: FT Whisper 0.100→0.094）、細部だけではほとんど得られなかった（0.100）。
一方、Zipformerは4~kHz未満のみの融合では悪化せず、悪化の原因はCE-SEが作る4~kHz以上の成分であった。
喉マイクに元々存在しない帯域に、Whisper向けに作られた成分が、Kaldi型の特徴量を用いる認識器には害となっている。
帯域ごとに重みを変えた12通りをdevで比較しても差は1\%前後で、$w=0.5$の一様な重みで十分であった。

\begin{table}[t]
  \centering
  \caption{融合の分解（test CER）。櫛除去: CE-SEの櫛状成分を除いてから融合、$<$4~kHz／$\geq$4~kHz: その帯域だけ融合}
  \label{tab:decomp}
  \small
  \setlength{\tabcolsep}{3.0pt}
  \begin{tabular}{lccccc}
    \toprule
    認識器 & 生成×4 & 融合 & 櫛除去 & $<$4~kHz & $\geq$4~kHz \\
    \midrule
    W-small* & 0.169 & 0.163 & 0.163 & 0.166 & 0.167 \\
    W-medium & 0.151 & 0.144 & 0.145 & 0.145 & 0.148 \\
    W-turbo & 0.131 & 0.125 & 0.126 & 0.126 & 0.128 \\
    W-FT & 0.100 & 0.092 & 0.093 & 0.094 & 0.099 \\
    Zipformer & 0.433 & 0.483 & 0.491 & \textbf{0.431} & 0.479 \\
    \bottomrule
  \end{tabular}
\end{table}

\subsection{他系統の認識器への害を減らす変更}
害はCE-SEがASR損失によって大きく変えた成分に由来すると考え、2つの簡単な変更を試した。
1つ目は重み空間での内挿\cite{wiseft}で、CE-SEの重みを初期値のTAPS SEの重みへ部分的に戻す。
\begin{equation}
  \theta(\alpha) = (1-\alpha)\,\theta_\mathrm{TAPS} + \alpha\,\theta_\mathrm{CE}
\end{equation}
$\alpha$はdevでWhisper-smallのCERのみを見て$\{0.1, 0.2, \ldots, 0.9\}$から選び、$\alpha=0.7$とした。
2つ目は、式(3)の振幅の平均を対数振幅の平均（幾何平均）に替える。

結果を表\ref{tab:wise}に示す。
内挿したSE（$\alpha=0.7$）は単体で、Whisper系6種すべてでCE-SEよりCERが低く（$-1.8$〜$-7.0$\%、base・medium・FT Whisper・Qwen3-ASRで有意）、CE-SEがほとんど効かなかったWhisper-baseでも$-7.0$\%（10/10話者）であった。
同時にXLS-R・Zipformerでの悪化も小さくなった（CE-SE比$-14.5$\%・$-21.8$\%）。
これを対数振幅で融合すると、Whisper系のCERは元の融合と$\pm1$\%以内で悪化はなく、Zipformerは0.483から0.442へ改善した（10/10話者）。
生成×4に対する悪化は11.7\%から2.1\%に縮み、XLS-Rも同等に戻った。
ただしdevのWhisper-smallでは融合の4通り（振幅・対数×CE-SE・内挿）の差は小さく（0.186〜0.189）、Whisper系での優位は示せない。

\begin{table}[t]
  \centering
  \caption{害を減らす変更（test CER）。融合A: 振幅の平均×CE-SE（表1の融合）、融合B: 対数振幅の平均×内挿したSE}
  \label{tab:wise}
  \small
  \setlength{\tabcolsep}{3.5pt}
  \begin{tabular}{lcccc}
    \toprule
    認識器 & CE-SE & 内挿$\alpha$=0.7 & 融合A & 融合B \\
    \midrule
    W-small* & 0.196 & 0.192 & 0.163 & 0.162 \\
    W-base & 0.267 & \textbf{0.248} & 0.201 & 0.199 \\
    W-medium & 0.169 & 0.165 & 0.144 & 0.143 \\
    W-turbo & 0.142 & 0.139 & 0.125 & 0.124 \\
    W-FT & 0.121 & 0.115 & 0.092 & 0.092 \\
    Qwen3 & 0.114 & 0.109 & 0.094 & 0.093 \\
    \midrule
    XLS-R & 0.279 & 0.238 & 0.163 & 0.161 \\
    Zipformer & 0.791 & 0.619 & 0.483 & \textbf{0.442} \\
    \bottomrule
  \end{tabular}
\end{table}

%────────────────────────────────────────
\section{考察}
結果は、喉マイク音声のWhisper系向けSEにおいて2つの成分が役割を分けていることを示す。
事前学習済み生成モデルは欠けた成分を自然な音声として補い、認識器の系統によらずCERを下げる。
CE-SEは特定の入力特徴量に合わせた変化を加え、その効果は損失に用いたWhisper-smallだけでなく、学習に用いていない他サイズのWhisperやQwen3-ASRにも及ぶ。
ただしWhisper-baseではCE-SE単体の改善が小さく（$-2.2$\%）、融合の上乗せもbaseとQwen3-ASRでは有意でなかった。Whisper系の中でも効き方には差がある。

この役割分担は、ASRの学習に用いていない表現での気導音声との距離にも表れる。
HuBERT-Large\cite{hubert}の全層で同一発話の気導音声とのフレームごとのコサイン距離を求めると、9認識器すべてで評価した17条件をまたいで、この距離は最も悪化した認識器のCER変化と強く相関した（Spearman $\rho=0.92$）。
一方、Whisper-smallのCERとの相関は弱く（$\rho=0.37$）、CE-SEは気導音声から離れる（距離0.202、TAPS SEは0.130）にもかかわらずWhisper系のCERを下げる。
融合は距離を生成×4の0.092から0.111へ増やす代わりに、Whisper系のCERを下げている。

生成モデル自体をCE損失で学習する試みが失敗したことは、Whisperに合わせる成分を生成過程に組み込むよりも、生成の外側で足すほうが容易であることを示唆する。
1段の見積もりに対する損失が多段の生成に持ち越されなかったのは、flow matchingの学習目標が出力の分布を合わせるものであり、個々の出力を特定の損失で最適化する構造になっていないためと考えられる。

重みを3割戻すだけでWhisper系の効果がむしろ増え、害が減ったことは、CE損失による学習がWhisper-smallに対して「効き過ぎ」ており、その行き過ぎた部分が他系統への害の多くを担っていることを示唆する（ただし追加学習を続けただけでもWhisper系のCERは1〜4\%下がっており、CE-SE自体の学習が十分でなかった可能性もある）。
実用上は、対数振幅×内挿の融合（融合B）がWhisper系の上乗せを保ちつつ他系統への害をほぼ避けられるため、最も扱いやすい。

本研究の限界として、計算量が大きい（4.3億パラメータのモデルを20段・4サンプル実行）こと、評価がTAPS（韓国語・1機種）のみであること、各モデルの学習が1回であること、Dissen型との比較が別の認識経路によることが挙げられる。
融合自体は既存の部品の組み合わせであり、今後は融合の出力を教師とした軽量なSEへの蒸留と、別言語・別機種の体内伝導マイクコーパス\cite{vibravox}での検証を行う。

%────────────────────────────────────────
\section{まとめ}
Whisper系の認識器に向けた喉マイクSEについて、事前学習済み生成モデルとASR損失で学習したSEの役割を調べた。
生成モデルの適応は9種の認識器すべてでCERを下げ、その効果は事前学習に由来した。
生成モデルの4サンプル平均とWhisper用SEの振幅融合は、学習に用いていないWhisper系にも上乗せをもたらし、TAPS SE比で26〜37\%CERを下げ、既存のASR損失によるSEより良かった。
一方、生成モデル自体をWhisperの損失で学習する方法は効果がなかった。
融合の上乗せは両帯域のスペクトル包絡から生じ、Whisper用の成分のうち4~kHz以上のものは他系統の認識器を悪化させた。
この害は、Whisper用SEの重みを初期値へ3割戻す内挿と対数振幅での融合により、Whisper系の上乗せを保ったままほぼ解消できた。

\footnotesize
\let\oldbibitem\bibitem
\renewcommand{\bibitem}{\setlength{\itemsep}{0pt}\setlength{\parsep}{0pt}\oldbibitem}
\bibliographystyle{IEEEtran}
\bibliography{refs}

\end{document}
